\documentclass[aspectratio=169]{beamer}

% Theme Settings
\usetheme{Madrid}
\usecolortheme{beaver} % Red/Gray theme matches typical Econ styling
\setbeamertemplate{navigation symbols}{}

% Packages
\usepackage{booktabs}
\usepackage{tikz}
\usepackage{amsmath}
\usepackage{graphicx}

% Title Information
\title{Unit 2: Supply and Demand}
\subtitle{Elasticity, Taxation, and Consumer Choice}
\author{Doğukan Günkördü}
\date{}

\begin{document}

%---------------------------------------------------------
% Title Slide
%---------------------------------------------------------
\begin{frame}
    \titlepage
\end{frame}

%---------------------------------------------------------
% Table of Contents
%---------------------------------------------------------
\begin{frame}{Learning Objectives}
    In this chapter, we will cover:
    \begin{itemize}
        \item Price elasticity of demand and supply
        \item Cross-price and Income elasticity
        \item Consumer and Producer Surplus
        \item Impact of Taxes (Deadweight Loss & Incidence)
        \item Price Ceilings and Floors
        \item Trade, Tariffs, and Quotas
        \item Consumer Choice Theory (Utility Maximization)
    \end{itemize}
\end{frame}

%---------------------------------------------------------
% SECTION 1: ELASTICITY OF DEMAND
%---------------------------------------------------------
\section{Elasticity of Demand}

\begin{frame}{Introduction to Elasticity}
    \textbf{The Law of Demand} states that consumers buy more when prices fall. 
    \vspace{0.5cm}
    
    However, \textbf{Elasticity} measures \textit{how much} more or less they buy.
    \begin{itemize}
        \item It measures \textbf{responsiveness}.
        \item \textbf{Inelastic:} Consumers are insensitive to price changes (quantity changes a little).
        \item \textbf{Elastic:} Consumers are sensitive to price changes (quantity changes a lot).
    \end{itemize}
\end{frame}

\begin{frame}{Calculating Price Elasticity of Demand ($E_d$)}
    \begin{block}{The Formula}
    \[ E_d = \left| \frac{\% \Delta \text{ Quantity Demanded}}{\% \Delta \text{ Price}} \right| \]
    \end{block}
    
    \textbf{Note:}
    \begin{itemize}
        \item Use absolute values (drop the negative sign).
        \item \textbf{Calculating \% Change:} 
        \[ \frac{\text{New} - \text{Old}}{\text{Old}} \]
        \textit{Mnemonic: "N - OOO" (New minus Old Over Old)}
    \end{itemize}
\end{frame}

\begin{frame}{Elasticity Coefficient Values}
    \begin{table}[]
    \centering
    \begin{tabular}{|l|c|l|}
    \hline
    \textbf{Type} & \textbf{Value} & \textbf{Description} \\ \hline
    Perfectly Inelastic & $= 0$ & "Mr. Stick" (Vertical Line) \\ \hline
    Relatively Inelastic & $< 1$ & Insensitive to price \\ \hline
    Unit Elastic & $= 1$ & \% $\Delta Q = \% \Delta P$ \\ \hline
    Relatively Elastic & $> 1$ & Sensitive to price \\ \hline
    Perfectly Elastic & $\infty$ & "Mr. Flat" (Horizontal Line) \\ \hline
    \end{tabular}
    \end{table}
    \vspace{0.2cm}
    \footnotesize{\textbf{Tip:} "Inelastic" starts with "I" (Vertical). "Elastic" has horizontal dashes like "E".}
\end{frame}

\begin{frame}{The Total Revenue Test}
    Total Revenue ($TR$) = Price $\times$ Quantity.
    
    \begin{itemize}
        \item \textbf{Elastic Demand ($E_d > 1$):} $P$ and $TR$ move in \textbf{opposite} directions.
        \begin{itemize}
            \item Price $\uparrow$ $\rightarrow$ $TR \downarrow$
        \end{itemize}
        \item \textbf{Inelastic Demand ($E_d < 1$):} $P$ and $TR$ move in the \textbf{same} direction.
        \begin{itemize}
            \item Price $\uparrow$ $\rightarrow$ $TR \uparrow$
        \end{itemize}
        \item \textbf{Unit Elastic ($E_d = 1$):} $TR$ remains unchanged when Price changes.
    \end{itemize}
\end{frame}

\begin{frame}{Elasticity Along a Linear Demand Curve}
    Elasticity varies along a downward-sloping linear demand curve.
    \vspace{0.5cm}
    \begin{itemize}
        \item \textbf{Top Left:} Elastic Range ($E_d > 1$)
        \item \textbf{Middle:} Unit Elastic ($E_d = 1$) — Total Revenue is maximized here.
        \item \textbf{Bottom Right:} Inelastic Range ($E_d < 1$)
    \end{itemize}
    \vspace{0.5cm}
    \textbf{Mnemonic:} \textit{EUI} (Elastic, Unit, Inelastic) $\rightarrow$ "Eat Up Idiots"
\end{frame}

\begin{frame}{Determinants of Elasticity}
    Three questions to determine if a good is Elastic or Inelastic:
    \begin{enumerate}
        \item \textbf{Are there adequate substitutes?}
        \begin{itemize}
            \item Yes $\rightarrow$ Elastic (e.g., Toyota cars).
            \item No $\rightarrow$ Inelastic (e.g., Insulin).
        \end{itemize}
        \item \textbf{Can the purchase be delayed?}
        \begin{itemize}
            \item Long time horizon $\rightarrow$ Elastic.
            \item Urgent/Emergency $\rightarrow$ Inelastic.
        \end{itemize}
        \item \textbf{Does it require a large \% of income?}
        \begin{itemize}
            \item Large \% (e.g., Boat) $\rightarrow$ Elastic.
            \item Small \% (e.g., Salt) $\rightarrow$ Inelastic.
        \end{itemize}
    \end{enumerate}
\end{frame}

%---------------------------------------------------------
% SECTION 2: OTHER ELASTICITIES
%---------------------------------------------------------
\section{Other Elasticities}

\begin{frame}{Cross-Price Elasticity of Demand (CPED)}
    Measures how a price change in Product Y affects Quantity of Product X.
    
    \[ CPED = \frac{\% \Delta \text{ Qty of Product X}}{\% \Delta \text{ Price of Product Y}} \]
    
    \begin{itemize}
        \item \textbf{Positive (+):} Substitute Goods (e.g., Price of Hot Dogs $\uparrow$, Demand for Burgers $\uparrow$).
        \item \textbf{Negative (-):} Complementary Goods (e.g., Price of Hot Dogs $\uparrow$, Demand for Buns $\downarrow$).
    \end{itemize}
\end{frame}

\begin{frame}{Income Elasticity of Demand}
    Measures how responsiveness of demand is to a change in income.
    
    \[ E_i = \frac{\% \Delta \text{ Quantity Demanded}}{\% \Delta \text{ Income}} \]
    
    \begin{itemize}
        \item \textbf{Positive (+):} Normal Good (Income $\uparrow$, Demand $\uparrow$).
        \item \textbf{Negative (-):} Inferior Good (Income $\uparrow$, Demand $\downarrow$).
    \end{itemize}
\end{frame}

\begin{frame}{Price Elasticity of Supply}
    Measures how sensitive producers are to changes in price.
    \[ E_s = \frac{\% \Delta \text{ Quantity Supplied}}{\% \Delta \text{ Price}} \]
    
    \textbf{Time is the key determinant:}
    \begin{itemize}
        \item \textbf{Short Run:} Supply is more inelastic (hard to build new factories quickly).
        \item \textbf{Long Run:} Supply is more elastic (firms can adjust production capacity).
    \end{itemize}
\end{frame}

%---------------------------------------------------------
% SECTION 3: SURPLUS AND PRICE CONTROLS
%---------------------------------------------------------
\section{Surplus \& Price Controls}

\begin{frame}{Consumer and Producer Surplus}
    \begin{columns}
        \column{0.5\textwidth}
        \textbf{Consumer Surplus (CS):}
        \begin{itemize}
            \item Difference between max price a consumer is willing to pay and actual price.
            \item Area \textit{below} Demand and \textit{above} Price.
        \end{itemize}
        
        \column{0.5\textwidth}
        \textbf{Producer Surplus (PS):}
        \begin{itemize}
            \item Difference between actual price and lowest price a seller is willing to accept.
            \item Area \textit{above} Supply and \textit{below} Price.
        \end{itemize}
    \end{columns}
    \vspace{0.5cm}
    \centering
    \textbf{Total Surplus = CS + PS} (Maximized at Equilibrium)
\end{frame}

\begin{frame}{Price Ceilings}
    \textbf{Price Ceiling:} Maximum legal price (e.g., Rent Control).
    \begin{itemize}
        \item Must be set \textbf{BELOW} equilibrium to be binding (effective).
        \item \textbf{Result:} Quantity Demanded > Quantity Supplied.
        \item Creates a \textbf{Shortage}.
        \item Creates \textbf{Deadweight Loss (DWL)}.
    \end{itemize}
\end{frame}

\begin{frame}{Price Floors}
    \textbf{Price Floor:} Minimum legal price (e.g., Minimum Wage, Ag. Supports).
    \begin{itemize}
        \item Must be set \textbf{ABOVE} equilibrium to be binding (effective).
        \item \textbf{Result:} Quantity Supplied > Quantity Demanded.
        \item Creates a \textbf{Surplus}.
        \item Creates \textbf{Deadweight Loss (DWL)}.
    \end{itemize}
\end{frame}

%---------------------------------------------------------
% SECTION 4: TAXATION
%---------------------------------------------------------
\section{Taxation}

\begin{frame}{Excise Taxes}
    An \textbf{Excise Tax} is a per-unit tax on production.
    \begin{itemize}
        \item Shifts the Supply curve vertically upward by the amount of the tax.
        \item Creates a "wedge" between the price consumers pay ($P_c$) and price producers receive ($P_s$).
        \item \textbf{Tax Revenue} = Tax Amount $\times$ Quantity Sold.
        \item \textbf{Deadweight Loss} = Loss of total surplus due to reduced quantity.
    \end{itemize}
\end{frame}

\begin{frame}{Tax Incidence (Who pays?)}
    The burden of the tax falls on the side of the market that is \textbf{more inelastic} (less responsive).
    
    \begin{table}[]
    \centering
    \begin{tabular}{|l|l|}
    \hline
    \textbf{Elasticity Comparison} & \textbf{Incidence} \\ \hline
    Demand is Inelastic & Consumers pay more \\ \hline
    Supply is Inelastic & Producers pay more \\ \hline
    Perfectly Inelastic Demand & Consumers pay \textbf{all} \\ \hline
    Perfectly Elastic Demand & Producers pay \textbf{all} \\ \hline
    \end{tabular}
    \end{table}
    \footnotesize{Mnemonic: If you can't leave the market (inelastic), you pay the tax.}
\end{frame}

%---------------------------------------------------------
% SECTION 5: TRADE
%---------------------------------------------------------
\section{Trade & Tariffs}

\begin{frame}{International Trade & Tariffs}
    \textbf{World Price ($P_w$):}
    \begin{itemize}
        \item If $P_w < P_{\text{Domestic}}$: Country imports the good. Consumer surplus increases; Producer surplus decreases.
    \end{itemize}

    \textbf{Tariffs (Tax on Imports):}
    \begin{itemize}
        \item Increases price above World Price (but usually below domestic autarky price).
        \item Reduces imports.
        \item Generates Tariff Revenue for government.
        \item Creates Deadweight Loss.
    \end{itemize}
\end{frame}

%---------------------------------------------------------
% SECTION 6: CONSUMER CHOICE
%---------------------------------------------------------
\section{Consumer Choice}

\begin{frame}{Utility Theory}
    \begin{itemize}
        \item \textbf{Total Utility (TU):} Total satisfaction from consuming a good.
        \item \textbf{Marginal Utility (MU):} Extra satisfaction from consuming \textbf{one more} unit.
    \end{itemize}
    
    \textbf{Diminishing Marginal Utility:}
    As you consume more of a good, the additional utility from each new unit decreases.
    \begin{center}
        \textit{"The first scoop of ice cream is great; the sixth scoop makes you sick."}
    \end{center}
\end{frame}

\begin{frame}{Utility Maximization Rule}
    Consumers have limited budgets. To maximize total utility, they should spend their budget such that:
    
    \[ \frac{MU_x}{P_x} = \frac{MU_y}{P_y} \]
    
    \textbf{The Logic:}
    \begin{itemize}
        \item If $\frac{MU_x}{P_x} > \frac{MU_y}{P_y}$, you get more "bang for your buck" with Good X. Buy more X, less Y.
        \item As you buy more X, $MU_x$ falls (diminishing returns) until the ratio equalizes.
    \end{itemize}
\end{frame}

\begin{frame}{Sample Utility Problem}
    \textbf{Goal:} Spend \$52. Souvenirs (\$8), Games (\$4).
    
    \begin{itemize}
        \item Step 1: Calculate $MU$ per dollar ($\frac{MU}{P}$) for every item.
        \item Step 2: Buy the item with the highest $\frac{MU}{P}$.
        \item Step 3: Repeat until budget is exhausted.
    \end{itemize}
    \vspace{0.5cm}
    \textit{This method ensures you are always getting the maximum utility for every dollar spent.}
\end{frame}

\begin{frame}{End of Unit 2 Review}
    \centering
    \Large \textbf{Questions?}
\end{frame}

\end{document}
